\documentclass[12pt, a4paper]{article}
\usepackage[utf8]{inputenc}
\usepackage[greek, english]{babel}
\usepackage{alphabeta}
\usepackage{amsmath, amssymb}
\usepackage{booktabs}
\usepackage{geometry}
\usepackage{float}

\geometry{a4paper, margin=1.2cm}

\title{\textbf{Απαντήσεις Εργασίας Κρυπτογραφίας}\\[1ex] \large Google Cybersecurity Seminars}
\author{Ζωίδης Βασίλειος}
\date{\today}

\begin{document}
\maketitle
\section*{Άσκηση 1 – Vigenère: Ανάκτηση Κλειδιού (Γνωστό Μήκος)}

Το κρυπτογράφημα που μας δίνεται είναι το \textbf{LXFOPVEFRNHR}. Γνωρίζουμε ότι το κλειδί έχει μήκος 5 χαρακτήρες και το κείμενο αποτελείται από πεζά αγγλικά γράμματα χωρίς σημεία στίξης.

\subsection*{α) Εύρεση του κλειδιού με ανάλυση ανά στήλη}

Εφόσον το μήκος του κλειδιού είναι 5, το $i$-οστό γράμμα του κρυπτογραφήματος έχει κρυπτογραφηθεί με το $(i \bmod 5)$-οστό γράμμα του κλειδιού. Ο τύπος αποκρυπτογράφησης είναι $P_i \equiv (C_i - K_i) \bmod{26}$. Χωρίζουμε το κρυπτογράφημα σε 5 στήλες ανάλογα με τη θέση του γράμματος (τύπου Caesar):
\begin{table}[htbp]
\begin{center}
\renewcommand{\arraystretch}{1.2}
\begin{tabular}{cc|c|c|c|c}
\textbf{mod5=} & \textbf{1} & \textbf{2} & \textbf{3} & \textbf{4} & \textbf{0} \\
\midrule
& L (11) & X (23) & F (5) & O (14) & P (15) \\
& V (21) & E (4)  & F (5) & R (17) & N (13) \\
& H (7)  & R (17) &       &        &        \\
\end{tabular}
\end{center}
\end{table}

\subsection*{β) Αποκρυπτογράφηση}
Δοκιμάζουμε τη λέξη-κλειδί \texttt{LEMON}, αφού αποτελεί κλασικό ακαδημαϊκό παράδειγμα στην κρυπτογραφία Vigenère. Αποκρυπτογραφώντας το μήνυμα βρίσκουμε:

\begin{center}
$K_{LEMON} = (11, 4, 12, 14, 13)$
\begin{align*}
P_1 &= (11 - 11)\bmod{26} = 0 \text{ (a)} \\
P_2 &= (23 - 4)\bmod{26} = 19 \text{ (t)} \\
P_3 &= (5 - 12)\bmod{26} = 19 \text{ (t)} \\
P_4 &= (14 - 14)\bmod{26} = 0 \text{ (a)} \\
P_5 &= (15 - 13)\bmod{26} = 2 \text{ (c)} \\
P_6 &= (21 - 11)\bmod{26} = 10 \text{ (k)} \\
P_7 &= (4 - 4)\bmod{26} = 17 \text{ (a)} \\
P_8 &= (5 - 12)\bmod{26} = 19 \text{ (t)} \\
P_9 &= (17 - 14)\bmod{26} = 3 \text{ (d)} \\
P_{10} &= (13 - 13)\bmod{26} = 0 \text{ (a)} \\
P_{11} &= (7 - 11)\bmod{26} = 22 \text{ (w)} \\
P_{12} &= (17 - 4)\bmod{26} = 13 \text{ (n)}
\end{align*}
\end{center}

Όπως βλέπουμε, το αποτέλεσμα σχηματίζει την πλήρως κατανοητή φράση \texttt{attackatdawn}. Συνεπώς, η υπόθεσή μας επιβεβαιώνεται απόλυτα.

\section*{Άσκηση 2 – RSA: Πλήρης Κύκλος με Ταχείες Υψώσεις}

Μας δίνονται οι πρώτοι αριθμοί $p=197$ και $q=211$.

\subsection*{α) Υπολογισμοί $n$, $\varphi(n)$ και $d$}

Υπολογίζουμε το $n$ πολλαπλασιάζοντας τα $p$ και $q$:
\[ n = p \times q = 197 \times 211 = 41567 \]

\noindent Υπολογίζουμε τη συνάρτηση Euler $\varphi(n)$:
\[ \varphi(n) = (p-1)(q-1) = 196 \times 210 = 41160 \]

\noindent Μας δίνεται το δημόσιο εκθετικό $e = 24377$. Πρέπει να βρούμε το ιδιωτικό εκθετικό $d$, τέτοιο ώστε $e \times d \equiv 1 \pmod{\varphi(n)}$. Για να το βρούμε, χρησιμοποιούμε τον Εκτεταμένο Αλγόριθμο του Ευκλείδη για τους αριθμούς 41160 και 24377. Η διαδικασία προς τα πίσω (Back-substitution) μας δίνει:
\[ 1 = 17393 \times 24377 - 10301 \times 41160 \]
Επομένως, το ιδιωτικό εκθετικό είναι:
\[ d = 17393 \]
\subsection*{β) Κρυπτογράφηση και Αποκρυπτογράφηση Μηνύματος ($m=1234$)}

\textbf{Κρυπτογράφηση:}\\
Χρησιμοποιούμε το δημόσιο κλειδί $(e, n) = (24377, 41567)$ και τον τύπο κρυπτογράφησης $C = M^e \bmod{n}$:
\[ C = 1234^{24377} \bmod{41567} \Rightarrow \]
\[ C = 26476 \]

\noindent \textbf{Αποκρυπτογράφηση:}\\
Για να επαληθεύσουμε το αποτέλεσμα, παίρνουμε το κρυπτογράφημα $C=26476$ και χρησιμοποιούμε το ιδιωτικό κλειδί $(d, n) = (17393, 41567)$ με τον τύπο αποκρυπτογράφησης $M = C^d \bmod{n}$:
\[ M = 26476^{17393} \bmod{41567} \Rightarrow \] 
\[ M = 1234 \texttt{ επαληθεύτηκε} \] 

\subsection*{γ) Αναλυτικά Βήματα Μεθόδου Square-and-Multiply}

Για την \textbf{κρυπτογράφηση} με τη μέθοδο \textbf{Square-and-Multiply}, μετατρέπουμε αρχικά τον εκθέτη $e=24377$ στο δυαδικό του ανάπτυγμα:
    \[24377_{10} = 101111100111001_{2}\]
\noindentΣαρώνοντας τα bits από αριστερά προς τα δεξιά και με αρχική τιμή $R = 1$, η διαδικασία διαμορφώνεται ως εξής:
\begin{itemize}
    \item Για κάθε bit εκτελείται \textbf{Square} (ύψωση του τρέχοντος αποτελέσματος στο τετράγωνο modulo $n$).
    \item Εάν το τρέχον bit είναι \textbf{$1$}, εκτελείται επιπλέον \textbf{Multiply} (πολλαπλασιασμός με τη βάση $m=1234\bmod{n}$).
\end{itemize}

\begin{table}[H]
\centering
\renewcommand{\arraystretch}{1.2}
\begin{tabular}{l l l}
\textbf{Bit} & \textbf{Λειτουργίες} & \textbf{Πράξεις και Υπόλοιπα (mod 41567)} \\
\midrule
\textbf{1} & \textbf{Square, Multiply} & $R = 1^2 \equiv 1$, $R = 1 \times 1234 \equiv 1234$ \\
\textbf{0} & \textbf{Square} & $R = 1234^2 \equiv 26344$ \\
\textbf{1} & \textbf{Square, Multiply} & $R = 26344^2 \equiv 3704$, $R = 3704 \times 1234 \equiv 39933$ \\
\textbf{1} & \textbf{Square, Multiply} & $R = 39933^2 \equiv 9668$, $R = 9668 \times 1234 \equiv 583$ \\
\textbf{1} & \textbf{Square, Multiply} & $R = 583^2 \equiv 7353$, $R = 7353 \times 1234 \equiv 11996$ \\
\textbf{1} & \textbf{Square, Multiply} & $R = 11996^2 \equiv 40629$, $R = 40629 \times 1234 \equiv 6384$ \\
\textbf{1} & \textbf{Square, Multiply} & $R = 6384^2 \equiv 19796$, $R = 19796 \times 1234 \equiv 28435$ \\
\textbf{0} & \textbf{Square} & $R = 28435^2 \equiv 29508$ \\
\textbf{0} & \textbf{Square} & $R = 29508^2 \equiv 18115$ \\
\textbf{1} & \textbf{Square, Multiply} & $R = 18115^2 \equiv 23327$, $R = 23327 \times 1234 \equiv 21154$ \\
\textbf{1} & \textbf{Square, Multiply} & $R = 21154^2 \equiv 22961$, $R = 22961 \times 1234 \equiv 26747$ \\
\textbf{1} & \textbf{Square, Multiply} & $R = 26747^2 \equiv 33939$, $R = 33939 \times 1234 \equiv 22757$ \\
\textbf{0} & \textbf{Square} & $R = 22757^2 \equiv 39363$ \\
\textbf{0} & \textbf{Square} & $R = 39363^2 \equiv 35844$ \\
\textbf{1} & \textbf{Square, Multiply} & $R = 35844^2 \equiv 39500$, $R = 39500 \times 1234 \equiv \mathbf{26476}$ \\
\end{tabular}
\end{table}

Επομένως, όντως το κρυπτογράφημα είναι \textbf{$C = 26476$}.

\section*{Άσκηση 3 – Διάδοση Σφάλματος στα Μπλοκ CBC και CFB}

Ο τρόπος που το σφάλμα στο ένα bit του $C_1$ διαδίδεται στην αποκρυπτογράφηση εξαρτάται απόλυτα από τον τρόπο λειτουργίας του εκάστοτε block cipher:

\subsection*{Αποκρυπτογράφηση με CBC (Cipher Block Chaining): $P_i = D_k(C_i) \oplus C_{i-1}$}

\begin{itemize}
    \item \textbf{Για το μπλοκ $P_1'$:} Το αλλοιωμένο $C_1$ εισέρχεται απευθείας στη συνάρτηση αποκρυπτογράφησης. Λόγω της ιδιότητας της διάχυσης (avalanche effect) που χαρακτηρίζει τους block ciphers, η έξοδος $D_k(C_1)$ τυχαιοποιείται πλήρως. Επομένως, το $P_1'$ θα καταστραφεί τελείως. Επηρεάζονται \textbf{όλα τα 128 bits} (με πιθανότητα περίπου 50\% για το καθένα να αλλάξει κατάσταση).
    \item \textbf{Για το μπλοκ $P_2'$:} Ο τύπος είναι $P_2 = D_k(C_2) \oplus C_1$. Εδώ το $C_2$ είναι σωστό, άρα η αποκρυπτογράφηση $D_k(C_2)$ είναι απολύτως σωστή. Στη συνέχεια, το αποτέλεσμα γίνεται XOR απευθείας με το εσφαλμένο $C_1$. Μια πράξη XOR με ένα μόνο λανθασμένο bit θα μεταβάλει ακριβώς το ίδιο αντίστοιχο bit στο αποτέλεσμα. Επομένως, επηρεάζεται \textbf{ακριβώς 1 bit} στο $P_2'$ (αυτό που βρίσκεται στην ίδια ακριβώς θέση με το αλλοιωμένο bit του $C_1$).
\end{itemize}

\subsection*{Αποκρυπτογράφηση με (Cipher Feedback): $P_i = C_i \oplus E_k(C_{i-1})$}

\begin{itemize}
    \item \textbf{Για το μπλοκ $P_1'$:} Ο τύπος για το πρώτο μπλοκ διαμορφώνεται ως $P_1 = C_1 \oplus E_k(IV)$. Το IV και η κρυπτογράφησή του παραμένουν σωστά. Επειδή το αλλοιωμένο $C_1$ μπαίνει απευθείας στην πράξη XOR, το σφάλμα μεταφέρεται άμεσα στο plaintext. Επομένως, επηρεάζεται \textbf{ακριβώς 1 bit} (το αντίστοιχο με το σφάλμα) στο $P_1'$.
    \item \textbf{Για το μπλοκ $P_2'$:} Ο τύπος είναι $P_2 = C_2 \oplus E_k(C_1)$. Εδώ το λανθασμένο $C_1$ εισέρχεται μέσα στη συνάρτηση κρυπτογράφησης $E_k$. Πάλι, εξαιτίας του avalanche effect, το αποτέλεσμα της $E_k(C_1)$ τυχαιοποιείται πλήρως. Όταν αυτό το εντελώς κατεστραμμένο μπλοκ γίνει XOR με το (κατά τα άλλα σωστό) $C_2$, θα παραγάγει ένα εξίσου κατεστραμμένο αποτέλεσμα. Επομένως, επηρεάζονται \textbf{όλα τα 128 bits} στο $P_2'$ (πάλι με περίπου 50\% πιθανότητα αλλαγής κατάστασης για το καθένα).
\end{itemize}
\end{document}